% SEGMENT OLP-0673-S01
% Part: proof-theory
% Chapter: normalization
% Section: normalization-thm

\documentclass[../../../include/open-logic-section]{subfiles}

\begin{document}

\olfileid{pt}{nor}{thm}

\olsection{இயல்பாக்கத் தேற்றம்}

!!^a{proof} இல் வெட்டுத் தொடர்கள் இல்லையெனில் அது
\emph{இயல்பு} வடிவில் உள்ளது. அதற்கு வரிசைமாற்று
உருமாற்றமோ குறைப்பு உருமாற்றமோ பயன்படுத்த முடியாது
என்பதற்கும் இதற்கும் சமானம் என்பது தெளிவு.
!!a{proof} ஐ \emph{இயல்பாக்குதல்} என்பது வெட்டுகள்
எதுவும் எஞ்சாதவரை வரிசைமாற்று மற்றும் குறைப்பு
உருமாற்றங்களைப் பயன்படுத்தி அதை இயல்பு வடிவ நிறுவலாக
மாற்றுவதாகும். இவ்வாறு எப்போதும் செய்யலாம் என்பதே
\emph{இயல்பாக்கத் தேற்றம்}.

\begin{thm}\ollabel{thm:normalization}
எந்த \Log{N1i}-!!{proof}~$\delta$ வையும் இயல்பாக்கலாம்;
அதாவது வெட்டுகளற்ற !!a{proof}~$\delta'$ ஆக மாற்றலாம்.
\end{thm}

% SEGMENT OLP-0673-S02
\begin{proof}
  $\delta$ இன் வெட்டு நிலை, வெட்டு நீளம் ஆகியவற்றின்
  மீது அகரவரிசைத் தொகுத்தறிதலைப் பயன்படுத்துகிறோம்.
  ஒரு !!{proof} இன் வெட்டு நிலை குறைவாகவோ, அதே வெட்டு
  நிலையில் அதன் வெட்டு நீளம் குறைவாகவோ இருந்தால் அதை
  இயல்பாக்கலாம் என்பதே தொகுத்தறிதல் கருதுகோள்.

  வெட்டு நீளம்~$0$ எனில், வெட்டுகள் இல்லை; $\delta$
  ஏற்கெனவே இயல்பு வடிவில் உள்ளது. எனவே $\delta$ இன்
  வெட்டு நீளம்~$> 0$ எனக் கொள்வோம்; அதன் வெட்டு நிலை
  பூச்சியமாகவும் இருக்கலாம். மிகுநிலை வெட்டுத்
  தொடர்களில் மிக மேலுள்ளவற்றுள் மிக வலப்புறமாக இருப்பதைத்
  தேர்ந்தெடுப்போம். அதன் நீளம் $>1$ எனில் வரிசைமாற்று
  உருமாற்றத்தைப் பயன்படுத்துவோம்; நீளம்~$1$ எனில் அதற்குரிய
  குறைப்பு உருமாற்றத்தைப் பயன்படுத்துவோம். இதனால் கிடைக்கும்
  புதிய !!{proof} இன் வெட்டு நிலை மாறாமல் வெட்டு நீளம்
  குறைந்திருக்கும், அல்லது வெட்டு நிலையே குறைந்திருக்கும்.
  இரு நேர்வுகளிலும் தொகுத்தறிதல் கருதுகோளைப் பயன்படுத்தலாம்.
\end{proof}

% SEGMENT OLP-0673-S03
மேலுள்ள நிறுவலில் பயன்படுத்திய உத்தி, மிக மேலுள்ள
மிக வலப்புறமான மிகுநிலை வெட்டை ஒவ்வொரு படியிலும்
கையாள்வது, !!a{proof} ஐ இயல்பு வடிவுக்கு மாற்றுவதற்கு
உருமாற்றங்களை ஒரு குறிப்பிட்ட வரிசையில் பயன்படுத்த வேண்டும்
என்று தோன்றச் செய்கிறது. வியப்பாக, அவற்றின் வரிசை
உண்மையில் பொருட்டல்ல. வரிசைமாற்று, குறைப்பு
உருமாற்றங்களை எந்த வரிசையில் பயன்படுத்தினாலும் இறுதியில்
இயல்பு வடிவ !!{proof} கிடைக்கும்.

% SEGMENT OLP-0673-S04
தொடர் மற்றும் வெட்டு ஆகியவற்றின் வரையறைகளை
\Log{N2i} க்கும் எளிதில் அமைக்கலாம். \Log{N1i} க்கான
வரிசைமாற்று, குறைப்பு உருமாற்றங்களும் நேரடியாக
\Log{N2i} க்கு மாற்றப்படுகின்றன. எனவே இயல்பாக்கத்
தேற்றம் \Log{N2i} க்கும் பொருந்துகிறது.

\begin{thm}\ollabel{thm:normalization-N2i}
எந்த \Log{N2i}-!!{proof}~$\delta$ வையும் இயல்பாக்கலாம்;
அதாவது வெட்டுகளற்ற !!a{proof}~$\delta'$ ஆக மாற்றலாம்.
\end{thm}

\end{document}
